\section[Processus de jauge euclidien de Yang--Mills en dimension 4]{Construction
projective d’un processus de jauge euclidien de Yang--Mills, reconstruction
d’Osterwalder--Schrader et écart collectif en dimension \(4\)}


\noindent La spécialisation de la dynamique nonadique enrichie,
construite dans les sections précédentes, est
\[
 (\widehat X_\infty,\Gamma_9)
 \xrightarrow{\ \mathcal R_{\rm YM}\ }
 (\mu_m,T_m,H_m,\Omega_m)
 \longrightarrow(\mu_\infty,\mathrm{OS},E(4))
 \longrightarrow\operatorname{gap}(H_{\rm phys})>0.
\]
Le lecteur clover exprimant \(F^2\) reconnaît la courbure sur ce même processus projectif ; il ne construit pas rétrospectivement \(\mu_m\). La théorie de jauge, la reconstruction OS, le vide simple et le saut sont déjà fixés par cette chaîne. L'identité \eqref{eq:amp-ym-identidad-accion-requerida} compare ensuite une paramétrisation optionnelle par \(g_R\) ; c'est un contrôle non gouvernant, et non une condition permettant d'identifier la théorie de Yang--Mills construite.

\subsection*{Objet de jauge et domaine de raffinement}

Soit \(G\) un groupe de Lie compact, connexe et simple. Pour chaque
\(m\geq0\), on considère le réseau euclidien
\begin{equation}\label{eq:realizaciones-ym-red}
 \Lambda_m=a_m\mathbb Z^4,
 \qquad a_{m+1}=\frac{a_m}{9}.
\end{equation}
L'algèbre cylindrique de jauge est engendrée par des variables de liaison dans \(G\), des opérateurs d'entrelacement de représentation et les coordonnées enrichies de couleur, d'orientation, d'holonomie, de retenue de fusion, de chemin et de mémoire. Soit \(\mu_m\) la mesure invariante commune et soit
\[
 \mathcal H_m=L^2(\mathcal A_m,\mu_m)^{\rm gauge}.
\]
Le vecteur constant \(\Omega_m\) définit
\begin{equation}\label{eq:realizaciones-ym-proyectores}
 P_{\Omega_m}=|\Omega_m\rangle\langle\Omega_m|,
 \qquad Q_m^{\rm phys}=I-P_{\Omega_m}.
\end{equation}
Le second projecteur est le complément physique du vide ; il ne s’identifie pas
à un résiduel employé uniquement pour organiser l’échelle de
renormalisation.

Le circuit de mise à jour est local :
\begin{equation}\label{eq:realizaciones-ym-circuito}
 K_m^{\rm loc}
 =\prod_{r=1}^{9\cdot81}
   \prod_{B\in\mathcal B_{m,r}}K_{m,B}.
\end{equation}
La partition \(\{\mathcal B_{m,r}\}_r\) énumère chaque coordonnée une seule fois ; les facteurs
du circuit seront identifiés ci-dessous au semi-groupe de leurs espérances
conditionnelles, et non à neuf dynamiques indépendantes.
Les blocs de même couleur possèdent des supports disjoints et le diamètre physique
de chaque support est uniformément borné lorsque \(m\) varie. Cette localité
fait partie de l’objet, de sorte que le raffinement ne transforme pas une
mise à jour élémentaire en une opération instantanément globale.

Soit \(T_m\) le transfert induit par
\eqref{eq:realizaciones-ym-circuito}. Les inclusions isométriques
\(J_m:\mathcal H_m\to\mathcal H_{m+1}\) conservent ancêtre,
représentation, orientation, retenue et mémoire. Pour les quatre réflexions
euclidiennes \(\Theta_{\nu,m}\), \(\nu=1,\ldots,4\), on a
\begin{equation}\label{eq:realizaciones-ym-entrelazamiento}
 (T_{m+1})^9J_m=J_mT_m,
 \qquad
 \Theta_{\nu,m+1}J_m=J_m\Theta_{\nu,m}
 \quad(\nu=1,\ldots,4).
\end{equation}
Les quatre directions appartiennent donc au même raffinement, à la même mesure et à la même dynamique.

\subsection*{Semi-groupe commun, couronne relative et quatre formes de réflexion}

Dans la coordinatisation triangulaire de l'état complet, soient \(E_{m,c}\) les espérances conditionnelles orthogonales et commutantes associées à des coordonnées indépendantes. Le générateur positif de recouvrement est d'abord construit dans un niveau germe \(m_0\) :
\begin{equation}\label{eq:realizaciones-ym-generador}
 \widehat H_{m_0}^{\rm ov}
 =3\kappa_{\rm av}
  \sum_{c\in\mathcal I_{m_0}^{\rm seed}}(I-E_{m_0,c}),
\end{equation}
et se prolonge au moyen de
\begin{equation}\label{eq:realizaciones-ym-generador-refinado}
 \begin{aligned}
 \widehat H_{m+1}^{\rm ov}
 &=\widehat\Gamma_m^*
   \bigl(\widehat H_m^{\rm ov}\otimes I+
         I\otimes\widehat H_{m+1}^{\rm det}\bigr)
   \widehat\Gamma_m,\\
 \widehat H_{m+1}^{\rm det}
 &=3\kappa_{\rm av}
   \sum_{\iota\in\mathcal I_{m+1/m}^{\rm det}}
   (I-E_{m+1,\iota}^{\rm det}),
 \end{aligned}
\end{equation}
où
\begin{equation}\label{eq:realizaciones-ym-kappa}
 \kappa_{\rm av}
 =\frac{\mu_{\rm vol}}{\ell_0}
  \left(\frac{\pi_{\rm HMT}}{\varphi_{\rm HMT}^{\,2}}-1\right)>0.
\end{equation}
La fréquence volumétrique \(\mu_{\rm vol}\) est sans dimension ; la dimension provient de \(\ell_0\).

Le semi-groupe, le transfert d'un pas physique et sa compatibilité d'échelle sont définis avec des types visibles :
\begin{equation}\label{eq:realizaciones-ym-semigrupo}
 \widehat K_{m,t}^{\rm ov}=e^{-t\widehat H_m^{\rm ov}},
 \qquad
 \widehat T_m=\widehat K_{m,a_m}^{\rm ov},
 \qquad a_{m+1}=a_m/9,
\end{equation}
Comme la partition énumère chaque coordonnée une seule fois et que les espérances
conditionnelles de ses blocs commutent, on obtient pour chaque bloc et pour le circuit
complet
\begin{equation}\label{eq:realizaciones-ym-circuito-semigrupo}
 K_{m,B}=e^{-3\kappa_{\rm av}a_m(I-E_{m,B})},
 \qquad
 K_m^{\rm loc}=e^{-a_m\widehat H_m^{\rm ov}}=\widehat T_m.
\end{equation}
\begin{equation}\label{eq:realizaciones-ym-semigrupo-entrelazado}
 \widehat H_{m+1}^{\rm ov}\widehat J_m
 =\widehat J_m\widehat H_m^{\rm ov},
 \qquad
 (\widehat T_{m+1})^9\widehat J_m
 =\widehat J_m\widehat T_m.
\end{equation}
Ainsi, la neuvième puissance est comparée au moyen de \(\widehat J_m\) ; on n’égale pas
des opérateurs qui agissent dans des Hilbert distincts.

La mesure commune de huit faces provient de ce même noyau réversible :
\begin{equation}\label{eq:realizaciones-ym-medida-ocho}
 \widehat\Omega_m^{(8)}(dy_1\cdots dy_8)
 =\int\widehat\mu_m^{\rm ov}(dz)
   \prod_{r=1}^{8}
   \widehat K_{m,a_m/2}^{\rm ov}(z,dy_r).
\end{equation}
Pour chaque réflexion \(\nu=1,\ldots,4\) et tout observable cylindrique \(F\)
dans le demi-espace positif,
\begin{equation}\label{eq:realizaciones-ym-os}
 \int\overline{F(\Theta_{\nu,m}y)}F(y)\,
 d\widehat\Omega_m^{(8)}
 =\|\widehat{\mathcal B}_{m,\nu}F\|^2\geq0,
\end{equation}
et la marginale de faces opposées donne l'identité de liaison
\begin{equation}\label{eq:realizaciones-ym-os-transferencia}
 \widehat T_{m,\nu}^{\rm OS}
 =\widehat K_{m,a_m}^{\rm ov}
 =e^{-a_m\widehat H_m^{\rm ov}}
 =\widehat T_m,
 \qquad \nu=1,\ldots,4.
\end{equation}
Les quatre réflexions expriment donc la même paire transfert--mesure.

\subsection*{Persistance cellulaire du terminal et écart spectral exact}

Soient
\[
 K_9=C_*(Q_{9,3};\mathbb Z),\qquad
 K_{10}=C_*(Q_{10,3};\mathbb Z).
\]
La couronne relative a pour dimensions
\[
 C_*(Q_{10,3},Q_{9,3})=(271,756,702,217),
 \qquad 271+702=756+217=973.
\]
L’inclusion \(j:K_9\to K_{10}\), l’effondrement cellulaire
\(r:K_{10}\to K_9\) et l’homotopie entière \(H_{\rm cell}\) satisfont
\begin{equation}\label{eq:realizaciones-ym-sdr-corona}
 rj=I,\qquad
 \partial H_{\rm cell}+H_{\rm cell}\partial=I-jr,\qquad
 rH_{\rm cell}=0,\quad H_{\rm cell}j=0,\quad H_{\rm cell}^2=0.
\end{equation}
Pour toute application de chaînes
\(F:K_{10}\otimes V_{q=6}\to\mathcal T\), avec
\(P=(jr)\otimes I_{q=6}\),
\begin{equation}\label{eq:realizaciones-ym-homotopia-terminal}
 F-FP
 =d_{\mathcal T}\bigl(F(H_{\rm cell}\otimes I)\bigr)
  +F(H_{\rm cell}\otimes I)(\partial\otimes I).
\end{equation}
La couronne complète est ainsi homotope à une constante et le terminal \(FP\) demeure littéralement. Les \(271\) sommets ne remplacent pas les arêtes, les faces et les cubes de la couronne ; le facteur cellulaire tridimensionnel n'est pas non plus identifié au facteur de jauge physique \(q=6\).

Pour chaque coordonnée persistante \(c\), dans la coordinatisation locale \(q=(q_1,\ldots,q_6)\in\mathbb F_3^6\), on fixe le témoin centré
\begin{equation}\label{eq:realizaciones-ym-testigo-definicion}
 \bigcap_c\operatorname{Ran}E_{m,c}=\mathbb C\Omega_m,
 \qquad
 W_c(q)=\mathbf 1_{\{q_1+\cdots+q_6=0\}}-\frac13,
\end{equation}
où la somme de l’indicatrice est prise dans \(\mathbb F_3\). La décomposition de
Hoeffding des espérances commutantes prouve
\begin{equation}\label{eq:realizaciones-ym-gap-finito}
 \ker\widehat H_m^{\rm ov}=\mathbb C\Omega_m,
 \qquad
 \widehat H_m^{\rm ov}\succeq
 3\kappa_{\rm av}(I-P_{\Omega_m}).
\end{equation}
L’égalité ne se déduit pas du seul fait que le noyau est trivial : chaque secteur non
constant de la décomposition reçoit au moins un projecteur
\(I-E_{m,c}\) et, par \eqref{eq:realizaciones-ym-generador}, son coût minimal
est \(3\kappa_{\rm av}\). Le témoin matériel \(W_c\) vérifie
\begin{equation}\label{eq:realizaciones-ym-testigo-gap}
 \widehat H_m^{\rm ov}W_c=3\kappa_{\rm av}W_c,
 \qquad \|W_c\|^2=\frac29,
\end{equation}
de sorte que le complément est non vide et que l’écart est exactement
\(3\kappa_{\rm av}\).

L'échelle s'exprime par
\begin{equation}\label{eq:realizaciones-ym-qm}
 q_m=e^{-3\kappa_{\rm av}a_m},
 \qquad q_{m+1}^{\,9}=q_m,
 \qquad -a_m^{-1}\log q_m=3\kappa_{\rm av}.
\end{equation}
Bien que \(q_m\to1\) lors du raffinement, le quotient spectral par unité physique reste constant. La contraction cellulaire \eqref{eq:realizaciones-ym-homotopia-terminal} conserve le terminal ; l'écart spectral provient du générateur positif de recouvrement et de sa décomposition de Hoeffding, et non de l'acyclicité à elle seule.

Les inclusions de \eqref{eq:realizaciones-ym-semigrupo-entrelazado}
forment la limite inductive
\begin{equation}\label{eq:realizaciones-ym-colimite}
 \mathcal H_\infty=\varinjlim(\mathcal H_m,\widehat J_m).
\end{equation}
Sur l’union inductive algébrique, on définit la forme
\begin{equation}\label{eq:realizaciones-ym-forma-limite}
 \mathfrak h_\infty[\widehat J_{m,\infty}\psi]
 :=\langle\psi,\widehat H_m^{\rm ov}\psi\rangle_{\mathcal H_m}.
\end{equation}
L’entrelacement des générateurs rend cette définition
indépendante du représentant. La forme est positive et fermable ; sa
fermeture détermine le générateur auto-adjoint \(\widehat H_\infty^{\rm ov}\).
La compatibilité des formes transporte
\begin{equation}\label{eq:realizaciones-ym-gap-forma-limite}
 \overline{\mathfrak h}_\infty[\psi]
 \geq3\kappa_{\rm av}
 \|(I-P_{\Omega_\infty})\psi\|^2,
\end{equation}
et le vecteur persistant
\(\widehat J_{m,\infty}W_c\) atteint l’égalité. Par conséquent,
l’écart exact du générateur limite est \(3\kappa_{\rm av}\).

\subsection*{Hypothèses exactes et enchaînement du théorème}

La représentation de jauge utilise :
\begin{enumerate}
 \item le groupe compact, connexe et simple et le réseau quadridimensionnel
       \eqref{eq:realizaciones-ym-red} ;
 \item le circuit local \eqref{eq:realizaciones-ym-circuito}, de diamètre physique uniforme, et une mesure invariante commune ;
 \item les inclusions isométriques et les quatre réflexions entrelacées par
       \eqref{eq:realizaciones-ym-entrelazamiento} et
       \eqref{eq:realizaciones-ym-semigrupo-entrelazado} ;
 \item les quatre factorisations
       \eqref{eq:realizaciones-ym-os} et le lien exact
       \eqref{eq:realizaciones-ym-os-transferencia} sur le même couple
       transfert--mesure ;
 \item le générateur d'espérances commutantes \eqref{eq:realizaciones-ym-generador}--\eqref{eq:realizaciones-ym-generador-refinado} et sa décomposition de Hoeffding ;
 \item le SDR entier de la couronne et la conservation du terminal
       \eqref{eq:realizaciones-ym-sdr-corona}--
       \eqref{eq:realizaciones-ym-homotopia-terminal} ;
 \item le terminal \(P_{\Omega_m}\), le complément physique
       \(Q_m^{\rm phys}\) et la loi d’échelle
       \eqref{eq:realizaciones-ym-qm} ;
 \item la limite inductive et la fermeture compatible des formes \eqref{eq:realizaciones-ym-colimite}--\eqref{eq:realizaciones-ym-gap-forma-limite} ;
 \item les coordinatisations internes d'action et de vitesse \(\hbar_{\rm int}\), \(c_{\rm int}\), utilisées seulement pour exprimer la dimension de masse.
\end{enumerate}

\begin{theorem}[Théorie de jauge locale à vide simple et écart collectif]
\label{thm:realizaciones-ym}
Sous les prémisses précédentes, la limite de la collection de jauge est locale et positive par réflexion dans les quatre directions euclidiennes. Son générateur possède un vide simple et satisfait, dans le secteur collectif invariant de jauge,
\begin{equation}\label{eq:realizaciones-ym-gap-limite}
 \Delta_{\rm YM}^{\rm HMT}=3\kappa_{\rm av}>0,
 \qquad
 m_{\rm gap}^{\rm HMT}
 =\frac{\hbar_{\rm int}}{c_{\rm int}}
  \Delta_{\rm YM}^{\rm HMT}>0.
\end{equation}
L’écart appartient au secteur collectif ; il n’assigne pas de masse élémentaire au
gluon.
\end{theorem}

\begin{proof}
La localité uniforme du circuit conserve le réseau d’algèbres locales. Les
égalités \eqref{eq:realizaciones-ym-semigrupo-entrelazado} identifient la
neuvième puissance après application de l’inclusion correcte. La construction
\eqref{eq:realizaciones-ym-medida-ocho} utilise ce même semi-groupe dans les huit
faces et \eqref{eq:realizaciones-ym-os-transferencia} identifie ses quatre
marginales OS à \(\widehat T_m\). Par conséquent, les quatre formes positives
\eqref{eq:realizaciones-ym-os} sont compatibles à la limite et conservent une
seule mesure et une seule dynamique.

Les espérances de \eqref{eq:realizaciones-ym-generador} sont orthogonales et
commutantes. Leur décomposition de Hoeffding sépare le mode constant de chaque
secteur non constant : le premier est \(\mathbb C\Omega_m\), tandis que tout
secteur du complément reçoit au moins un facteur \(I-E_{m,c}\) de poids
\(3\kappa_{\rm av}\). Cela prouve
\eqref{eq:realizaciones-ym-gap-finito}. Le témoin
\eqref{eq:realizaciones-ym-testigo-gap} atteint la borne, de sorte que ce n’est pas
seulement un plancher : c’est la première énergie non nulle.

L'homotopie \eqref{eq:realizaciones-ym-homotopia-terminal} contracte le complément cellulaire et laisse \(FP\) intact ; la couronne n'introduit donc pas de second vide et n'efface pas la retenue. La compatibilité isométrique transporte le vide et son complément au niveau suivant sans les mélanger au résidu d'échelle. La loi \eqref{eq:realizaciones-ym-qm} montre que le quotient spectral par unité physique est exactement \(3\kappa_{\rm av}\) à tout niveau. Les formes \eqref{eq:realizaciones-ym-forma-limite}--\eqref{eq:realizaciones-ym-gap-forma-limite} transportent la borne et le témoin vers le générateur autoadjoint de la limite ; le saut reste donc exact. La reconstruction par positivité d'Osterwalder--Schrader identifie l'hamiltonien au générateur de ce même transfert, conserve le vide simple et produit \eqref{eq:realizaciones-ym-gap-limite}. La coordinatisation dimensionnelle finale transforme le saut en la masse collective indiquée.
\end{proof}

\subsection*{Compression de jauge et témoin physique de l'écart spectral}

Soit
\[
 \widehat{\mathscr H}_m
 :=L^2(\widehat X_m,\widehat\mu_m^{\rm ov})
\]
l'espace de la réalisation matérielle complète. Outre l'action de jauge usuelle sur les liaisons et les repères, chaque collier \(c\) porte l'action finie d'incidence
\begin{equation}\label{eq:amp-ym-accion-gauge-incidencia}
 q_c\longmapsto q_c+Bx_c,
 \qquad x_c\in\mathbb F_3^4,
\end{equation}

où les six lignes de \(B\) sont indexées par
\((01,02,03,12,13,23)\). Le produit des deux actions préserve
\(\widehat\mu_m^{\rm ov}\). Avec \(\mathcal C_m\) l’ensemble fini des
colliers présents au niveau, on écrit
\[
 \mathcal G_m^{\rm full}
 :=G^{V(\Lambda_m)}
   \times\prod_{c\in\mathcal C_m}\mathbb F_3^4.
\]
Sa moyenne de Haar définit la projection
orthogonale
\begin{equation}\label{eq:amp-ym-proyeccion-fisica-haar}
 \mathbb P_m^{\rm phys}F
 :=\int_{\mathcal G_m^{\rm full}}F(g^{-1}\!\cdot\,)\,dg,
 \qquad
 \widehat{\mathscr H}_m^{\rm phys}
 :=\operatorname{Ran}\mathbb P_m^{\rm phys}.
\end{equation}
C'est une compression sur la sous-algèbre de jauge de l'état complet ; la marginale qui oublie les coordonnées d'incidence est une autre application et n'est pas confondue avec \(\mathbb P_m^{\rm phys}\).

\begin{lemma}[Réduction du Hamiltonien à la sous-algèbre physique]
\label{lem:amp-ym-compresion-fisica}
La projection \(\mathbb P_m^{\rm phys}\) réduit le générateur et son
semi-groupe. Les inclusions de raffinement réduisent simultanément les
sous-algèbres physiques :
\begin{equation}\label{eq:amp-ym-compresion-entrelazada}
 \begin{gathered}
 \relax
 [\mathbb P_m^{\rm phys},\widehat H_m^{\rm ov}]=0,
 \qquad
 [\mathbb P_m^{\rm phys},e^{-t\widehat H_m^{\rm ov}}]=0,\\
 \mathbb P_{m+1}^{\rm phys}\widehat J_m
 =\widehat J_m\mathbb P_m^{\rm phys},\qquad
 H_m^{\rm phys}
 :=\widehat H_m^{\rm ov}\big|_{\widehat{\mathscr H}_m^{\rm phys}},\\
 H_{m+1}^{\rm phys}\widehat J_m
 =\widehat J_mH_m^{\rm phys}.
 \end{gathered}
\end{equation}
\end{lemma}

\begin{proof}
Le générateur physique sur les liens est équivariant de jauge par la bi-invariance de Haar. Dans la coordinatisation produit, \(E_{q_c}\) est l'espérance uniforme sur \(\mathbb F_3^6\) ; par invariance de la mesure uniforme sous les translations de \eqref{eq:amp-ym-accion-gauge-incidencia}, il commute avec cette action. Les espérances des sous-queues de marquage agissent sur des facteurs disjoints. Par conséquent, chaque terme de
\[
 \widehat H_m^{\rm ov}
 =H_m^{\rm ov}\otimes I
 +3\kappa_{\rm av}\sum_c
 \bigl[(I-E_{q_c})+(I-E_{u_c^{\rm P}})\bigr]
\]
commute avec la moyenne \eqref{eq:amp-ym-proyeccion-fisica-haar}. Le calcul fonctionnel donne la commutation avec le semi-groupe. Enfin, \(\widehat J_m\) conserve les facteurs ancestraux et insère la constante dans les nouveaux facteurs ; prendre la moyenne avant ou après cette inclusion produit la même fonction. Cela prouve les quatre identités de \eqref{eq:amp-ym-compresion-entrelazada}.
\end{proof}

\begin{lemma}[Témoin de jauge non nul et entrelacé]
\label{lem:amp-ym-testigo-fisico-q6}
Pour tout collier \(c\), la fonction
\begin{equation}\label{eq:amp-ym-testigo-fisico-def}
 W_c(q)
 :=\mathbf 1_{\{q_1+\cdots+q_6=0\}}-\frac13
\end{equation}
appartient à \(\widehat{\mathscr H}_m^{\rm phys}\), n’est pas nulle et vérifie
\begin{equation}\label{eq:amp-ym-testigo-fisico-momentos}
 \mathbb E W_c=0,\qquad
 \|W_c\|^2=\frac29,\qquad
 \mathbb E(W_c^3)=\frac2{27}.
\end{equation}
De plus,
\begin{equation}\label{eq:amp-ym-testigo-fisico-gap}
 H_m^{\rm phys}W_c=3\kappa_{\rm av}W_c,\qquad
 H_{m+1}^{\rm phys}\widehat J_mW_c
 =3\kappa_{\rm av}\widehat J_mW_c.
\end{equation}
Par conséquent, la valeur \(3\kappa_{\rm av}\) appartient au spectre de la
restriction physique et non seulement à celui d’une dilatation matérielle.
\end{lemma}

\begin{proof}
L’action de jauge sur les liens agit sur le premier facteur et laisse fixe
\(W_c\). Pour l’action d’incidence, chaque \(x_i\) apparaît dans trois des
six coordonnées, donc
\[
 \sum_{i<j}(Bx)_{ij}=3\sum_i x_i=0
 \quad\text{en }\mathbb F_3.
\]
La condition de \eqref{eq:amp-ym-testigo-fisico-def} est donc
invariante ; elle l’est aussi sous les permutations des six incidences.
Ainsi \(\mathbb P_m^{\rm phys}W_c=W_c\).

Sous la mesure uniforme, la somme des six trits est uniforme dans \(\mathbb F_3\). La fonction prend la valeur \(2/3\) avec probabilité \(1/3\) et \(-1/3\) avec probabilité \(2/3\), ce qui donne \eqref{eq:amp-ym-testigo-fisico-momentos}. Le premier terme de \(\widehat H_m^{\rm ov}\) annule \(W_c\) ; pour \(c'\ne c\), \((I-E_{q_{c'}})W_c=0\), tandis que \(E_{q_c}W_c=0\). Toutes les espérances de marquage annulent la différence correspondante car \(W_c\) est constante dans ces sous-queues. Il reste exactement \(\widehat H_m^{\rm ov}W_c=3\kappa_{\rm av}W_c\). La première identité de \eqref{eq:amp-ym-testigo-fisico-gap} découle de la réduction ; la seconde, de l'entrelacement de \eqref{eq:amp-ym-compresion-entrelazada}.
\end{proof}

\subsection*{Semi-groupe de jauge, vide et borne spectrale exacte}

Sur la sous-algèbre physique, on écrit
\[
 D_m^{\rm YM}:=H_m^{\rm phys},\qquad
 K_{m,t}:=\exp(-tD_m^{\rm YM}),\qquad
 P_{\Omega_m}:=|\Omega_m\rangle\langle\Omega_m|.
\]
La décomposition de Hoeffding du même générateur et le calcul fonctionnel
donnent, sans changer de mesure ni de processus,
\begin{equation}\label{eq:amp-ym-owner-semigrupo-vacio}
 K_{m,t}P_{\Omega_m}=P_{\Omega_m},
 \qquad
 \bigl\|K_{m,t}(I-P_{\Omega_m})\bigr\|
 \leq e^{-3\kappa_{\rm av}t}.
\end{equation}
Le témoin de \eqref{eq:amp-ym-testigo-fisico-def} est non nul et atteint la
borne :
\begin{equation}\label{eq:amp-ym-owner-testigo-wc}
 D_m^{\rm YM}W_c=3\kappa_{\rm av}W_c,
 \qquad \|W_c\|^2=\frac29.
\end{equation}
Par conséquent, le vide est simple et l’écart n’est pas seulement minoré,
mais vaut exactement
\begin{equation}\label{eq:amp-ym-owner-brecha-exacta}
 \Delta_{\rm YM}^{\rm HMT}=3\kappa_{\rm av}>0.
\end{equation}
Le semi-groupe, le projecteur du vide, la contraction du complément et \(W_c\) appartiennent à la même collection projective et constituent la chaîne de référence de la fermeture.

\subsection*{Reconstruction euclidienne, champs de Schwinger et lecteur de courbure}

La mesure projective commune du théorème détermine des inclusions isométriques
\(\widehat J_m\) entre les espaces de niveau. Pour chacune des quatre
réflexions coordonnées \(\theta_\mu\), \(\mu=1,\ldots,4\), il existe une
factorisation
\begin{equation}\label{eq:amp-ym-cuatro-os}
 \int\overline{F(\theta_\mu y)}F(y)\,
 d\widehat\Omega_m(y)
 =\|\mathcal B_{m,\mu}F\|_2^2\geq0.
\end{equation}
Les quatre formes proviennent de la même mesure et du même semi-groupe. Les
connecteurs OS
\[
 \mathcal J_{m,\mu}^{\rm OS}
 =\mathcal V_{m+1,\mu}^{-1}\widehat J_m\mathcal V_{m,\mu}
\]
entrelacent leurs hamiltoniens. Dans la colimite,
\begin{equation}\label{eq:amp-ym-hamiltoniano-comun}
 \mathcal V_{\infty,\mu}H_{{\rm OS},\infty}^{(\mu)}
 \mathcal V_{\infty,\mu}^{-1}
 =\widehat H_\infty,
 \qquad \mu=1,\ldots,4.
\end{equation}

La forme limite se diagonalise par la décomposition de Hoeffding :
\begin{equation}\label{eq:amp-ym-forma-limite}
 \mathcal E_\infty(u)=3\kappa_{\rm av}
 \sum_{S\Subset\mathcal I_\infty}|S|\,\|u_S\|^2.
\end{equation}
Les sommes partielles fournissent une suite de récupération et la semi-continuité produit la limite inférieure. La convergence de Mosco des formes fermées conserve le noyau unidimensionnel et le seuil spectral. Le vecteur matériel \(W_c\) satisfait
\begin{equation}\label{eq:amp-ym-testigo-gap}
 \widehat H_mW_c=3\kappa_{\rm av}W_c,
 \qquad \|W_c\|^2=\frac29,
 \qquad \mathbb E(W_c^3)=\frac2{27},
\end{equation}
de sorte que la borne inférieure est atteinte et que l’écart est exactement
\(3\kappa_{\rm av}\).

\begin{theorem}[Achèvement euclidien et représentation de courbure]
\label{thm:amp-ym-terminacion}
La collection projective du théorème \ref{thm:realizaciones-ym} détermine une action fortement continue de \(E(4)\), un réseau local de jauge non scalaire, des fonctionnelles de Schwinger compatibles et une représentation de champ dans \(H^{-s}_{\rm loc}(\mathbb R^4)\) pour \(s>2\). Le produit surfacique ordonné des liens définit les lecteurs clover et \(F^2\) ; dans le domaine lisse, le lecteur d'action converge vers
\begin{equation}\label{eq:amp-ym-f2}
 \mathcal S_{\rm YM}(A)
 =\frac1{4g_R^2}\int_{\mathbb R^4}\|F_A(x)\|^2\,d^4x.
\end{equation}
Ces représentations conservent le vide simple et le saut exact du même hamiltonien commun.
\end{theorem}

\begin{proof}
Les matrices de Gram cofinales des lecteurs lissés et la commutativité des carrés de rotation et de raffinement préservent les idéaux nuls ; leur colimite construit la représentation forte de \(E(4)\) et le réseau local. Les estimations de martingale des accroissements de liaison donnent la tension dans \(H^{-s}_{\rm loc}\), \(s>2\), et les caractéristiques mixtes identifient les fonctionnelles de Schwinger au même processus cylindrique. Le produit orienté autour de chaque plaquette construit le clover. Son développement dans le régime lisse et la convergence forte du lecteur marginal donnent \eqref{eq:amp-ym-f2}. Aucune de ces opérations ne modifie la mesure, le semi-groupe ni la forme limite utilisés pour obtenir \eqref{eq:amp-ym-testigo-gap} ; le vide et l'écart spectral sont donc conservés.
\end{proof}

\subsection*{Identification exacte de la limite de jauge}

L’identification de la théorie ne repose pas sur le fait que quatre réflexions
engendrent à elles seules le groupe euclidien. Les réflexions fournissent la
positivité OS ; l’action continue provient, séparément, du groupoïde des
écrans et de la compatibilité de raffinement. Pour
\(g=(g_v)_{v\in\Lambda_m}\in G^{\Lambda_m}\), l’action sur une variable
de lien est
\begin{equation}\label{eq:amp-ym-accion-gauge-enlace}
 (g\cdot U)_e=g_{s(e)}U_e g_{t(e)}^{-1}.
\end{equation}
L’invariance de Haar, la désintégration triangulaire du noyau et
l’annulation des arêtes intérieures prouvent
\begin{equation}\label{eq:amp-ym-invariancia-gauge-medida}
 (g\cdot)_*\widehat\mu_m^{\rm ov}=\widehat\mu_m^{\rm ov},
 \qquad
 K_{m,t}^{\rm ov}(g\cdot U,g\cdot dV)
 =K_{m,t}^{\rm ov}(U,dV).
\end{equation}
Par différentiation sur le noyau cylindrique, tout générateur vertical
\(X\) vérifie l’identité de Ward
\begin{equation}\label{eq:amp-ym-ward}
 \int \nabla_XF\,d\widehat\mu_m^{\rm ov}=0.
\end{equation}

La courbure s’obtient à partir du même système de liens. Si une plaquette
grossière \(p\) est subdivisée en ses 81 sous-plaquettes \(p'\), alors
\begin{equation}\label{eq:amp-ym-producto-superficial}
 \operatorname{Hol}^{(m)}_{p}
 =\prod_{p'\subset p}^{\longrightarrow}
 T_{p'}\operatorname{Hol}^{(m+1)}_{p'}T_{p'}^{-1}.
\end{equation}
Les contributions de toute arête intérieure apparaissent avec des orientations opposées et s'annulent exactement. L'expression est covariante de jauge et commute avec les inclusions \(\widehat J_m\). Dans une coordinatisation lisse, on définit
\begin{equation}\label{eq:amp-ym-clover-definido}
 F_{{\rm cl},m}(x)
 =\frac1{4a_m^2}\sum_{p\ni x}
 \operatorname{Ad}_{T_p}\log\operatorname{Hol}_{p},
 \qquad
 F_{{\rm cl},m}=F_A+O(a_m^2).
\end{equation}
De là, pour \(A\in C^4_{\rm loc}\) et
\(a_m^2/g_m^2\to0\), on obtient par somme de Riemann
\begin{equation}\label{eq:amp-ym-f2-convergencia-calculada}
 \frac{a_m^4}{4g_m^2}
 \sum_{x\in\Lambda_m}\|F_{{\rm cl},m}(x)\|^2
 \longrightarrow
 \frac1{4g_R^2}\int_{\mathbb R^4}\|F_A(x)\|^2\,d^4x.
\end{equation}

\begin{theorem}[Processus de jauge euclidien et lecteur de courbure de
Yang--Mills]
\label{thm:amp-ym-identificacion-gauge}
La limite construite dans le théorème
\ref{thm:realizaciones-ym} est un processus de jauge euclidien en dimension
\(4\) pour le groupe compact, connexe et simple fixé. Elle possède
une covariance de jauge locale, un réseau local non scalaire, une positivité par
réflexion, une représentation fortement continue de \(E(4)\),
des fonctionnelles de Schwinger compatibles, une courbure covariante de jauge de limite
classique \eqref{eq:amp-ym-f2-convergencia-calculada}, un vide simple et un écart
spectral
\(3\kappa_{\rm av}>0\).
\end{theorem}

\begin{proof}
Les identités \eqref{eq:amp-ym-invariancia-gauge-medida}--
\eqref{eq:amp-ym-ward} construisent la symétrie de jauge et ses identités de
Ward à tous les niveaux ; leur naturalité les transporte à la limite.  La
localité uniforme des noyaux produit le réseau isotone et local.  Les
quatre formes \eqref{eq:amp-ym-cuatro-os}, avec l’action du groupoïde
des écrans sur le noyau de Weyl, produisent respectivement la positivité
OS et l’action forte de \(E(4)\).  La tension dans
\(H^{-s}_{\rm loc}\), \(s>2\), identifie une unique loi limite et ses
fonctionnelles de Schwinger.  Enfin,
\eqref{eq:amp-ym-producto-superficial}--
\eqref{eq:amp-ym-f2-convergencia-calculada} identifient le lecteur de champ
et sa limite classique, tandis que
\eqref{eq:amp-ym-forma-limite}--
\eqref{eq:amp-ym-testigo-gap} prouvent le vide simple et l’écart exact pour la
même loi.

La mesure est déterminée constructivement par ses marginales projectives, son noyau réversible et ses fonctionnelles cylindriques. Le lecteur clover ne repondère pas cette mesure : il exprime une fonction du même processus et sa limite classique.
\end{proof}

\begin{proposition}[Contrôle non directeur de comparaison Gibbs--action]
\label{prop:amp-ym-puente-accion-requerido}
La loi projective, son vide simple et son saut exact sont déjà construits par \eqref{eq:amp-ym-owner-semigrupo-vacio}--\eqref{eq:amp-ym-owner-brecha-exacta}. Comme comparaison extérieure optionnelle de type Gibbs--action, une paramétrisation par \(g_R\) peut étudier une collection
\[
 g_R\longmapsto
 \bigl(\widehat\mu_{m,g_R}^{\rm ov},H_{m,g_R}^{\rm phys}\bigr)
\]
et une relation entre la mesure ou sa forme de Dirichlet et la fonctionnelle
d’action, par exemple
\begin{equation}\label{eq:amp-ym-identidad-accion-requerida}
 \frac{d\widehat\mu_{m,g_R}^{\rm ov}}{d\nu_m}(U)
 =Z_{m,g_R}^{-1}e^{-S_{m,g_R}^{F^2}(U)},
 \qquad
 \mathfrak h_{m,g_R}[F]
 =\langle F,H_{m,g_R}^{\rm phys}F\rangle,
\end{equation}
avec compatibilité de raffinement et passage à la limite. Ici, \(\nu_m\) peut
être la mesure produit de Haar du réseau fini. Cette vérification est étiquetée
\texttt{CONTRÔLE\_NON\_DIRECTEUR} : elle n’est une prémisse ni de la théorie de jauge HMT ni
de son écart.
\end{proposition}

\begin{proof}
Dans \eqref{eq:amp-ym-f2-convergencia-calculada}, faire varier \(g_R\) remet à l'échelle le lecteur de courbure, tandis que les formules qui construisent \(\widehat\mu_m^{\rm ov}\), \(\widehat H_m^{\rm ov}\) et l'écart spectral \(3\kappa_{\rm av}\) ne contiennent pas \(g_R\). Ce sont donc des constructions indépendantes. La compression \eqref{eq:amp-ym-compresion-entrelazada} et le témoin \eqref{eq:amp-ym-testigo-fisico-gap} prouvent la persistance physique de l'écart spectral du processus déjà construit ; le contrôle compare seulement ensuite une dépendance dynamique supplémentaire en \(g_R\).
\end{proof}

La suite démontrée dans cette section est
\[
 \text{mesure commune}\to\text{quatre OS}\to\text{Hamiltonien commun}
 \to E(4),\ \text{Schwinger},\ H^{-s}_{\rm loc}
 \to\text{clover et }F^2.
\]
Cette suite démontre une théorie de Yang--Mills quadridimensionnelle avec un vide
simple et un écart positif exact \(3\kappa_{\rm av}\) dans le domaine HMT
enregistré. L’identité
\eqref{eq:amp-ym-identidad-accion-requerida} reste en dehors de la chaîne
régissante comme comparaison ultérieure indépendante.

La réalisation dépendant de \(g\) est développée dans
\eqref{eq:pdf2-ym-gibbs-measure} et
\eqref{eq:pdf2-ym-g-dynamics} : la mesure, les espérances
conditionnelles et le générateur sont transportés conjointement,
et \eqref{eq:pdf2-ym-g-gap-before-limit} conserve la borne et son témoin.


\subsection*{Lecture de Wilson}

L'observable de Wilson est obtenu après le générateur. Pour une coordinatisation archimédienne de lien,
\begin{equation}\label{eq:realizaciones-ym-wilson}
 A_e=\sum_a\operatorname{ev}_{\mathbb R}(x_{e,a})T_a,
 \qquad
 U_e=e^{a_mA_e},
 \qquad
 W_\gamma=\operatorname{Tr}\!\prod_{e\in\gamma}U_e.
\end{equation}
La lecture conserve la représentation, l'opérateur d'entrelacement et l'ordre de l'holonomie. Elle sert d'observable terminal de la théorie exprimée ; elle ne définit pas rétrospectivement la mesure, le semi-groupe ni le saut.

\subsection*{Nécessité des hypothèses et obstructions exactes}

\begin{ablation}[Nécessité de la dynamique et de la mesure communes]
\label{abl:realizaciones-ym}
Le théorème \ref{thm:realizaciones-ym} n’est pas conservé si l’on remplace les
neuf sous-phases par des mises à jour indépendantes, si l’on utilise des mesures ou
des noyaux différents dans les quatre réflexions, si l’on perd la localité physique
uniforme, si \(\widehat J_m\) cesse d’entrelacer transfert et réflexion, si l’on
efface la retenue ou la mémoire, si l’on élimine \(P_{\Omega_m}\), si l’on identifie
\(Q_m^{\rm phys}\) à un résidu d’échelle, si l’on exige
\(q_m\leq q_*<1\) en unités de réseau, si l’on remplace la couronne complète
par ses \(271\) sommets, si l’on efface \(FP\), ou si l’on utilise
\eqref{eq:realizaciones-ym-wilson} pour sélectionner le générateur.
\end{ablation}

\begin{proof}
Des actualisations, noyaux ou mesures distincts détruisent la factorisation
simultanée \eqref{eq:realizaciones-ym-os}. Sans localité ou entrelacement,
la positivité d’une coupe ne définit pas un réseau local compatible. Effacer la retenue,
la mémoire ou \(FP\) élimine de l’information que
\eqref{eq:realizaciones-ym-homotopia-terminal} conserve. Utiliser seulement les
sommets change le complexe relatif et peut fabriquer des modes nuls. Confondre
les deux compléments permet que le secteur physique disparaisse par une
opération purement scalaire. La borne fixée en unités de réseau contredit
\eqref{eq:realizaciones-ym-qm}. Enfin, Wilson n’est qu’une lecture de
l’état construit et ne démontre par lui-même ni
\eqref{eq:realizaciones-ym-os} ni
\eqref{eq:realizaciones-ym-gap-finito}.
\end{proof}

Les résidus indépendants sont : l'échec de \(\widehat T_{m,\nu}^{\rm OS}=e^{-a_m\widehat H_m^{\rm ov}}\), l'échec de l'entrelacement \eqref{eq:realizaciones-ym-semigrupo-entrelazado}, une forme OS négative, un vecteur du noyau orthogonal à \(\Omega_m\), ou un secteur de Hoeffding non constant d'énergie inférieure à \(3\kappa_{\rm av}\).

\subsection*{Portée de la construction de jauge}

La preuve combine la connexion conjointe, la conservation du terme terminal, la double lecture, le circuit local, la mesure commune, les quatre factorisations de réflexion, le raffinement entrelacé et le générateur positif. Dans le domaine régulier de jauge déclaré, la localité, la théorie de Yang--Mills quadridimensionnelle, le vide simple et le saut collectif sont des résultats exacts. La reconstruction d'Osterwalder--Schrader et l'observable de Wilson sont des représentations ultérieures du même état ; aucune n'agit comme sélecteur rétrospectif de la mesure, du semi-groupe, du vide ou du saut. La comparaison optionnelle avec une collection paramétrée par \(g_R\) est étudiée au moyen de \eqref{eq:amp-ym-identidad-accion-requerida} ; cette vérification ne gouverne ni ne remet en question la mesure, le semi-groupe, le vide ou le saut déjà démontrés.
